\documentclass[11pt, a4paper, listof=numbered, captions=tableheading, headinclude, table, xcdraw]{scrreprt}

\usepackage[utf8]{inputenc}
\usepackage{tikz}
\usepackage{epigraph}
\usepackage{amsmath}
\usepackage{textgreek}
\usepackage[LGRgreek]{mathastext}
\usepackage{array}
\usepackage{color}
\usepackage{graphicx}
\usepackage{float}
\usepackage{helvet}
\usepackage{scrhack}
%\usepackage{circuitikz}
\renewcommand{\familydefault}{\sfdefault}

\setlength{\emergencystretch}{\hsize}
%\tolerance=9999
 
\newcounter{divide}

% für referenziertes Inhaltsverzeichnis
\usepackage[hidelinks, hypertexnames=false]{hyperref}

\usepackage{chngcntr}
\counterwithin*{chapter}{divide}

\renewcommand\epigraphflush{flushright}
\renewcommand\epigraphsize{\normalsize}
\setlength\epigraphwidth{0.7\textwidth}

\definecolor{titlepagecolor1}{cmyk}{1,.39,.0,.27}
\definecolor{titlepagecolor}{cmyk}{0,.82,.78,.30}

\DeclareFixedFont{\titlefont}{T1}{ppl}{b}{it}{0.5in}

\setlength{\parindent}{0pt}
\setlength{\parskip}{\baselineskip}

% Seitenränder festlegen
\usepackage[top=1.5cm,bottom=1.5cm, includeheadfoot]{geometry}


% ---------------------------------------
% Titeldefinition
\newcommand{\Titel}{Tech-Note\\\vspace{.5mm}SynPLC}
\newcommand{\Caption}{An OpenSource PLC system with STM32}
\newcommand{\SignateAuthor}{\textit{Made with \LaTeX}}
\newcommand{\Stand}{\textit{Stand: \today}}
% Sonstiges in titel.ltx und titel-meta.ltx
% ---------------------------------------

% Kopf- und Fußzeile
\usepackage[automark]{scrlayer-scrpage}



% --------------------------------------------------------------------
% Titelblatt
% --------------------------------------------------------------------

\makeatletter                       
\def\printauthor{         
    {\large \@author}}              
\makeatother
\author{
    Marco Müller \\ 
    \vspace*{12pt}
    www.marco-bastelt.de \\
    \texttt{marco@marco-bastelt.de}
    }

\renewcommand\epigraphflush{flushright}
\renewcommand\epigraphsize{\normalsize}
\setlength\epigraphwidth{0.7\textwidth}

\definecolor{Green}{gray}{0.55}

\newcommand\anglei{-45}
\newcommand\angleii{-45}
\newcommand\angleiii{-225}
\newcommand\angleiv{135}

\newcommand\titlepagedecoration{%
\begin{tikzpicture}[remember picture,overlay,shorten >= -10pt]

\coordinate (aux1) at ([yshift=-15pt]current page.north east);
\coordinate (aux2) at ([yshift=-410pt]current page.north east);
\coordinate (aux3) at ([xshift=-4.5cm]current page.north east);
\coordinate (aux4) at ([yshift=-150pt]current page.north east);

\begin{scope}[titlepagecolor!40,line width=12pt,rounded corners=12pt]
\draw
  (aux1) -- coordinate (a)
  ++(225:5) --
  ++(-45:5.1) coordinate (b);
\draw[shorten <= -10pt]
  (aux3) --
  (a) --
  (aux1);
\draw[opacity=0.6,titlepagecolor,shorten <= -10pt]
  (b) --
  ++(225:2.2) --
  ++(-45:2.2);
\end{scope}
\draw[titlepagecolor,line width=8pt,rounded corners=8pt,shorten <= -10pt]
  (aux4) --
  ++(225:0.8) --
  ++(-45:0.8);
\begin{scope}[titlepagecolor!70,line width=6pt,rounded corners=8pt]
\draw[shorten <= -10pt]
  (aux2) --
  ++(225:3) coordinate[pos=0.45] (c) --
  ++(-45:3.1);
\draw
  (aux2) --
  (c) --
  ++(135:2.5) --
  ++(45:2.5) --
  ++(-45:2.5) coordinate[pos=0.3] (d);   
\draw 
  (d) -- +(45:1);
\end{scope}
\end{tikzpicture}%
}
\newcommand\chapterdecoration{%
\begin{tikzpicture}[remember picture,overlay,shorten >= -10pt]
\coordinate (aux1) at ([yshift=-15pt]current page.north west);
\coordinate (aux2) at ([yshift=-310pt]current page.north west);
\coordinate (aux3) at ([xshift=4.5cm]current page.north west);
\coordinate (aux4) at ([yshift=-150pt]current page.north west);
\coordinate (aux5) at ([yshift=-535pt]current page.north west);
\coordinate (aux6) at ([yshift=-660pt]current page.north west);
\coordinate (aux7) at ([yshift=-630pt]current page.north west);
\coordinate (aux8) at ([xshift=-16.5cm]current page.south east);
\renewcommand\anglei{-135}
\renewcommand\angleii{135}
\renewcommand\angleiii{-45}
\renewcommand\angleiv{45}



\begin{scope}[titlepagecolor1!70,line width=6pt,rounded corners=8pt]
\draw[shorten <= -10pt]
  (aux2) --
  ++(\angleiii:3) coordinate[pos=0.4] (d) --
  ++(\anglei:3.1);

\draw
  (d) --
  ++(\angleiii:.1) --
 ++(\angleiii:1.7) --
  ++(\anglei:1.5) --
++(\angleiii:2.5) --
++(\anglei:2.5) --
++(\angleii:2) --
++(\angleiii:1.3) --
++(\anglei:1.3) 
 
coordinate[pos=0.3] (c); 
\begin{scope}[titlepagecolor1!40,line width=12pt,rounded corners=12pt]
\draw
  (aux7) -- coordinate (a)
  ++(\angleiii:5) --
  ++(\anglei:5.1) coordinate (b);
\draw[shorten <= -10pt]
  
  (b) -- 
++(\angleiv:2.5) --
++(\angleiii:5) 
  (aux8);
\draw[opacity=0.6,titlepagecolor1,shorten <= -10pt]
  (aux5) --
  ++(\angleiii:2.2) --
  ++(\anglei:2.2);
\draw[titlepagecolor1,line width=8pt,rounded corners=8pt,shorten <= -10pt]
  (aux6) --
  ++(\angleiii:0.8) --
  ++(\anglei:0.8);
\end{scope}
  
\end{scope}
\end{tikzpicture}%
}

% Old Code for reference
%
%
%
%\newcommand\titlepagedecoration{
%\begin{tikzpicture}[remember picture,overlay,shorten >= -10pt]
%
%\coordinate (aux1) at ([yshift=-15pt]current page.north east);
%\coordinate (aux2) at ([yshift=-410pt]current page.north east);
%\coordinate (aux3) at ([xshift=-4.5cm]current page.north east);
%\coordinate (aux4) at ([yshift=-150pt]current page.north east);
%
%\begin{scope}[titlepagecolor!40,line width=12pt,rounded corners=12pt]
%\draw
%  (aux1) -- coordinate (a)
%  ++(225:5) --
%  ++(-45:5.1) coordinate (b);
%\draw[shorten <= -10pt]
%  (aux3) --
%  (a) --
%  (aux1);
%\draw[opacity=0.6,titlepagecolor,shorten <= -10pt]
%  (b) --
%  ++(225:2.2) --
%  ++(-45:2.2);
%\end{scope}
%\draw[titlepagecolor,line width=8pt,rounded corners=8pt,shorten <= -10pt]
%  (aux4) --
%  ++(225:0.8) --
%  ++(-45:0.8);
%\begin{scope}[titlepagecolor!70,line width=6pt,rounded corners=8pt]
%\draw[shorten <= -10pt]
%  (aux2) --
%  ++(225:3) coordinate[pos=0.45] (c) --
%  ++(-45:3.1);
%\draw
%  (aux2) --
%  (c) --
%  ++(135:2.5) --
%  ++(45:2.5) --
%  ++(-45:2.5) coordinate[pos=0.3] (d);   
%\draw 
%  (d) -- +(45:1);
%\end{scope}
%\end{tikzpicture}
%}

\begin{document}
\begin{titlepage}

\vspace*{2cm}
\noindent
\titlefont \Titel\par
\epigraph{\Caption}
{\SignateAuthor\\ \Stand}
\vspace*{2cm}
\centering{\includegraphics[width=0.5\textwidth]{pictures/SynPLC.png}}
\null\vfill
\vspace*{1cm}
\noindent
\hfill
\begin{minipage}{0.4\linewidth}
    \begin{flushright}
        \printauthor
    \end{flushright}
\end{minipage}
%
\begin{minipage}{0.02\linewidth}
    \rule{1pt}{80pt}
\end{minipage}
\titlepagedecoration
\chapterdecoration
\end{titlepage}



% --------------------------------------------------------------------
% Verzeichnisse
% --------------------------------------------------------------------

\pagenumbering{Roman}
\renewcommand{\thechapter}{\Roman{chapter}}

\setcounter{tocdepth}{2} %Ausblenden von Subsections etc..
\cohead[I. Table of Contents]{I. Table of Contents}
\pagestyle{scrheadings}
\tableofcontents 	\newpage
\cohead[II. List of Tables]{II. List of Tables}
\pagestyle{scrheadings}
\listoffigures 		\newpage
\cohead[III. List of Figures]{III. List of Figures}
\pagestyle{scrheadings}
\listoftables 		\newpage
\pagenumbering{arabic}

\stepcounter{divide}

\renewcommand{\thechapter}{\arabic{chapter}}

% Kopfzeilen erst nach den Verzeichnissen

\lohead[\today]{\today}
\cohead[SynPLC V2.0]{SynPLC V2.0}
\rohead[Marco Müller]{Marco Müller}
\pagestyle{scrheadings}
 
% --------------------------------------------------------------------
% Ab hier der Inhalt
% --------------------------------------------------------------------

\chapter{About}

This document will describe the workings of the SynPLC.

The SynPLC is meant to be an OpenSource PLC system for home automation purposes. It can supply digital and analog IO as well as different modules for a variety of different applications, e.g. Phase-Fired Control of AC voltages. 

\section{The Origins}

This project was created of the idea that in theory one could build a LOGO!-type PLC system from Arduinos. After some research and a few failed attempts it evolved into the current state, using STM32 Microcontrollers. 


\section{License}
All resources and documents for the SynPLC project are distributed under the Creative Commons CC-BY-SA 4.0 license.

 				\newpage 	% Vorwort
\chapter{Summary}

The SynPLC project is a PLC system using STM32 microcontrollers and off-the-shelf hardware. For the current version the following modules are available:

\begin{itemize}
    \item CPU module
    \item Digital Output module
    \item Digital Input module
    \item Analog Output module
    \item Analog Input module
    \item Triac PFC module
\end{itemize}

The following modules are planned, but not yet under development:

\begin{itemize}
    \item PWM module
    \item Relais module
    \item Thermocouple module
\end{itemize}

For further information and technical specs, please look into the section  "Modules".

				\newpage
\chapter{Busses}

Every module has at least one common bus connector connecting via ribbon cable to the other modules. Alongside the 5V and GND voltage supply there also is the communication bus lines. On the modules there are two different busses possible: RS485 and CAN. 

In order to maintain a stable communication, consider the use of twisted-pair ribbon cables.

Both are working on a custom protocol, which is easy to implement on other controllers too. 

\section{RS485}

\subsection{Baudrate and Transmission speed}
The RS485 bus is transceiving at 250kBaud. At a maximum of 10 Bytes per frame, this results in a minimum rate of around 3000 frames per second.

\subsection{Frame Structure}
The frame structure of the protocol is as follows:
\vspace{0.5cm}

\begin{tabular}{cl}
    \textbf{Byte Number}     & \textbf{Data}                      \\
    1               & Startbyte (0x5A)          \\
    2               & Slave Adress              \\
    3               & Instruction and Register  \\
    4               & Data (Not used when Instruction == Scan or Read) \\
    5               & Data (only used at 12- or 16bit data blocks) \\
    6               & Checksum \\
    \\
    7               & Answer from Slave (ACK or NACK with Error code) \\
    8               & Answer from Slave (data, LSB at 12- or 16bit data blocks) \\
    9               & Answer from Slave (data, MSB at 12- or 16bit data blocks) \\
    10              & Answer from Slave (Checksum)
\end{tabular}
\vspace{0.5cm}

As you can see, most of the frame consists of the CPU talking to the slaves. Startbyte and Checksums are in place to make sure that an error free communication is possible. 

The Checksum itself is easily calculated. This is done by XOR-ing all sent bytes with the constant value 0xC5. In summary, the equation is as follows:

\begin{equation}
    Checksum = 0xC5 \oplus Byte\_1 \oplus Byte\_2 \oplus Byte\_3 (\oplus Byte\_4 \oplus Byte\_5)
\end{equation}

Equally, the slaves are calculating their checksum with the same method, XOR-ing all their bytes with 0xC5. 

If all the bytes are XOR-ed with each other and with the checksum itself and 0xC5 is still left, then the communication was error free and the data is uncompromized.

\subsection{Slave Addresses}
The slave addresses are an 8bit number which consist of a 3-bit module identifier and a 5-bit address settable via dip-switches on the respective module.

The module identifiers are definied as follows:

\begin{table}[H]
    \centering
        \begin{tabular}{|c|c|}
            \hline
            Identifier & Modul \\ \hline \hline
            000 & Digital Input \\ \hline
            001 & Digital Output \\ \hline
            010 & Analog Input \\ \hline
            011 & Analog Output \\ \hline
            100 & PWM Modul \\ \hline
            101 & Thyristor Modul \\ \hline
            110 & Relais Modul \\ \hline
            111 & \textit{reserved}    \\ \hline
        \end{tabular}
\end{table}

\subsection{Instructions}

In Byte 3 of the communication protocol are the instruction and the corresponding register specified. For more information on the registers, please refer to the corresponing module section in this document.

The instructions are at the 4 MSBs of the byte. They are as follows:

\begin{table}[H]
    \centering
        \begin{tabular}{|c|c|}
            \hline
            Instruction-Code & Instruction \\ \hline \hline
            0000 xxxx   & Bus Scanning (is slave available?) - no data byte \\ \hline
            0001 (reg)  & Read 8bit-Register\\ \hline
            0010 (reg)  & Write 8bit-Register \\ \hline
            0101 (reg)  & Read 16bit-Register \\ \hline
            0110 (reg)  & Write 16bit-Register \\ \hline
        \end{tabular}
\end{table}

\subsection{ACK, NACK and Error Codes}
In order to know weather the communication was successful, there is a byte used specifically for ACK/NACK and Error codes. This byte is divided into 2 nibbles, where the first nibble is the ACK/NACK and the second nibble is the error code. If everything is fine, the ACK nibble will be 0xC0 (1100). If there is an error present, it will display the NACK status with 0x03 (0011).

The following nibble will show which exactly went wrong:

\begin{table}[H]
    \centering
        \begin{tabular}{|c|c|}
            \hline
            Error Code & Meaning \\ \hline \hline
            0000 & No Error \\ \hline
            0001 & Target Register invalid \\ \hline
            0010 & Instruction invalid \\ \hline
            1100 & Processing Error - read Error Register \\ \hline
            1101 & Checksum Error - Send again \\ \hline
            1111 & System Error - read Error Register \\ \hline
        \end{tabular}
\end{table}				\newpage	% Bus Erklärung
\chapter{Modules V1}
All V1 modules are now discontinued. The successors are the V2 modules.

\vspace{2cm}

\textbf{Further information:}

V1 was meant to be a Proof of Concept system using easy-to-program ATmega328P with Arduino framework. 

The Hardware was not as refined as wanted and a lot of compromises were made resulting in high production costs and hard-to-source parts. 

\chapter{Modules V2}
These are the current module revisions. If not stated otherwise, every module is powered by a STM32F103 microcontroller. 

Every module is capable of both CAN and RS485, but not simultaneously. To select one of the two bus systems, set the Jumpers on the boards accodingly.

\vspace{1cm}

V2 was designed with the mindset of more generic circuitry without losing accuracy or stability. Even tho not every module was improved (talking about you, DO8...) major steps were made on the analog modules. In addition, the triac module (XP) is in testing now.

Also all V1 modules got discontinued and decommissioned.\pagebreak

\section{CPU}
The CPU is the heart of the PLC system and necessary for normal operation. It connects to all IO-modules and processes all the data.

The CPU has the following interfaces:

\begin{itemize}
    \item RS485 and CAN Businterface
    \item USB-Interface for serial monitor
    \item ESP01-Modul as Webinterface (read-only as of now)
    \item SD-card for structured text (WIP) and log files
    \item DC/DC-Converter as power supply for the whole system
\end{itemize}

The DC/DC-Converter uses an ESD-safe 24V DC-input to generate 5V DC at 10W output power. This voltage is available at the bus connector and providing power to all modules in the system.

It has no seperate ID since it is a one-of-a-kind module. In later revisions, there will be measures for redundancy available, but not for now.

\section{DO8}
The DO8 (Digital Output 8-channel) is providing 8 channels of 24V compatible outputs with a current rating of 0.5A each. The channels are galvanically isolated from the system voltage via optoisolators. 

The 24V Output voltage has to be supplied on the connector. 

To prevent major faults caused by short circuits, the output channels are measured back and compared to their set value. If a channel differs from its value, the module sets an error bit and shuts off the faulty channel.

On the front are 8 LEDs, indicating the status of the output channels. 

%hier noch die Register beschreiben

\section{DI}
The DI8 (Digital Input 8-channel) is providing 8 channels of 24V compatiple inputs, divided into two sets of 4 Inputs where both sets are galvanically isolated from each other. The maximum input current is limited to 2mA. The channels are galvanically isolated from the system voltage via optoisolators.

For normal operation, the system ground fromt the inputs need to be connected. 

The channels are ESD safe and can tolerate up to 33V DC. Reverse-Polarity and overcurrent protections are in place. 

A digital Low-Level is defined as 0..5 VDC.

A digital HIGH-Level is defined as 10..33 VDC.

% Hier die Register beschreiben

\section{AO}
The AO4 (Analog Output 4-channel) is providing 4 channels of analog output, each capable of supplying either 0..10V CV or 0..20mA CC at a resolution of 12bits. The switch between CV and CC is managed by a relay.

The output resolution is set at 16bits, giving a 0.15mV or 0.30$\mu$A resolution. 

The channels are electrically isolated from the system, so at least 15V with common analog grond have to be provided.

On the front are 4 LEDs indicating if the channel is actively being used. 

% Hier die Register beschreiben

\section{AI}
The AI4 (Analog Input 4-channel) is providing 4 channels of analog inputs, capable of meauring 0..10V or 0..20mA each. 

The input resolution is set at 12bits. At a maximum input volatage of 10V or input current of 20mA, this is giving a resolution of 2.4mV or 4.8mA.

The channels are electrically isolated from the system.


\section{XT}
The XT4 (Interface Triac 4-channel) is providing 4 triacs on two seperate voltage domains. The channels can be used to either switch AC voltage to a supply or even using them in a PFC (phase-fired controller) mode. 

For PFC mode each voltage domain has a Zero-Crossing-Detection (ZCD) and can control the triacs with PWM. The ZCD also calculates the frequency of the AC voltage automatically and compensates for frequency drift. 

Each channel is rated for up to 230V and 1A. Keep in mind that the triacs can become quite hot. 

All channels are electrically isolated from the system. The two voltage domains are also isolated from each other.  			\newpage	% Module V1



	

% --------------------------------------------------------------------
% abschließende Verzeichnisse
% --------------------------------------------------------------------

\pagestyle{empty}
\vspace*{\fill}
\begin{center}
\hrulefill\\
\vspace{4mm}
	\begin{HUGE}
		\textbf{Appendix}
	\end{HUGE}\\
\hrulefill
\end{center}
\vspace*{\fill}
\newpage
\pagestyle{plain}


\appendix

% --------------------------------------------------------------------
\end{document}
